% Research fragment. The general edge-half problem remains open here.
\section{Partial results for the cubic edge-half problem}
\label{sec:cubic-edge-frequency}

Write $c_e(G)$ for the number of cycles containing an edge $e$.
The edge-half problem asks for an edge with $c_e(G)\le c(G)/2$.
The two usual exceptions among connected simple cubic graphs are $K_4$ and $K_{3,3}$.
The results below settle several structural cases, but do not settle the general problem.
No priority claim is made for these elementary reductions.

\begin{proposition}\label{prop:cubic-small-structures}
Let $G$ be a connected simple cubic graph other than $K_4$ or $K_{3,3}$.
If $G$ has a bridge, a two-edge cut, a triangle, or a four-cycle, then some edge satisfies
\[
 c_e(G)\le\frac12c(G).
\]
\end{proposition}

\begin{proof}
A bridge lies in no cycle.
We may therefore suppose that $G$ is bridgeless.

First suppose that $\delta(S)$ is a two-edge cut.
Both induced shores $N_1,N_2$ are connected.
Otherwise one shore component has a cut of size at most one.
On each shore, the two cut endpoints are distinct.
If they coincided, the third edge at that endpoint would be a bridge.
The terminals have degree two within their shore.

Let $p_i$ count terminal paths in $N_i$, and put $c_i=c(N_i)$.
Every cycle crossing the cut uses a terminal path on each shore, so
\[
 c(G)=c_1+c_2+p_1p_2.
\]
Relabel the shores so that $c_1\le c_2$.
At one terminal of $N_1$, let $f,g$ be its two internal edges.
The terminal paths partition according to their first edge.
Choose $f$ to occur in at most $p_1/2$ terminal paths.
At most $c_1$ cycles wholly inside $N_1$ contain $f$.
Consequently,
\[
 c_f(G)\le c_1+\frac12p_1p_2\le\frac12c(G).
\]

Next suppose that $W$ is a triangle.
Put $c_0=c(G-V(W))$.
Let $S$ count paths in $G-E(W)$ joining distinct vertices of $W$.
Their internal vertices lie outside $W$, since every vertex of $W$ has outside degree one.
Every cycle meeting both $W$ and its exterior is one such path plus one of the two corresponding triangle arcs.
Thus
\[
 c(G)=c_0+1+2S,\qquad c_e(G)=1+S\quad(e\in E(W)).
\]
If $n=|V(G)|\ge6$, the exterior has $n-3$ vertices and $3n/2-6$ edges.
It therefore contains a cycle, since its edge count is at least its vertex count.
Hence $c_0\ge1$, and every triangle edge satisfies the required inequality.
The only simple cubic graph of order four is $K_4$.

Finally suppose that $W$ is a four-cycle.
A chord gives a triangle, so assume that $W$ is induced.
We claim the exact identity
\begin{equation}\label{eq:cubic-square-frequency}
 \sum_{e\in E(W)}c_e(G)=2c(G)+2-2c(G-V(W)).
\end{equation}
To prove it, group cycles by their edge sets outside $E(W)$.
A cycle with two attachment vertices has one outside path and two possible rim completions.
The two completed cycles use four rim edges in total.
A cycle with four attachment vertices has two disjoint outside paths.
Every connected completion uses a perfect matching of the four-cycle, and therefore uses exactly two rim edges.
The cycle $W$ itself uses four rim edges.
A cycle avoiding $E(W)$ avoids $V(W)$, since each rim vertex has outside degree one.
These cases prove~\eqref{eq:cubic-square-frequency}.

If $n\ge8$, the exterior has $n-4$ vertices and $3n/2-8$ edges.
Its edge count is at least its vertex count, so it contains a cycle.
Equation~\eqref{eq:cubic-square-frequency} gives
\[
 \min_{e\in E(W)}c_e(G)
 \le\frac14\sum_{e\in E(W)}c_e(G)\le\frac12c(G).
\]
For $n=6$, the complement of $G$ is a two-regular simple graph.
It is either a six-cycle or two triangles.
The corresponding cubic graphs are the triangular prism and $K_{3,3}$.
The former has a triangle and the latter is excluded.
\end{proof}

Thus any unresolved connected simple cubic example is three-edge-connected and has girth at least five.
The square identity also gives the quantitative bound
\[
 \min_{e\in E(W)}c_e(G)
 \le\frac12c(G)-\frac12\bigl(c(G-V(W))-1\bigr).
\]

\subsection{A weighted consequence of the four-regular classification}

\begin{proposition}\label{prop:cubic-matching-weight}
Let $G$ be a connected simple cubic graph on $n\ge6$ vertices.
For every perfect matching $M$ of $G$,
\begin{equation}\label{eq:cubic-matching-weight}
 \sum_{C\text{ a cycle of }G}2^{|E(C)\cap M|}\le2^n-\frac n2.
\end{equation}
If $G$ is bridgeless, then
\begin{equation}\label{eq:cubic-length-weight}
 \sum_{C\text{ a cycle of }G}2^{|E(C)|/3}\le2^n-\frac n2.
\end{equation}
\end{proposition}

\begin{proof}
Double every edge of $M$, obtaining a loopless four-regular multigraph $Q$.
Its cycles of length two are exactly the $n/2$ doubled matching edges.
Each longer cycle projects to one cycle $C$ of $G$ and has $2^{|E(C)\cap M|}$ lifts.
Hence
\[
 c(Q)=\frac n2+\sum_C2^{|E(C)\cap M|}.
\]
The four-regular classification gives $c(Q)\le2^n$.
Indeed, its underlying graph is cubic, so $Q$ is neither a doubled cycle nor $K_5$.
The remaining exceptional graph has order four.
This proves~\eqref{eq:cubic-matching-weight}.

Suppose now that $G$ is bridgeless.
Every odd cut has odd size and has at least three edges.
The vector assigning weight $1/3$ to every edge satisfies the perfect matching polytope constraints.
Edmonds' theorem therefore gives a probability distribution on perfect matchings with edge marginals $1/3$.
For a fixed cycle $C$, convexity gives
\[
 \mathbb E_M 2^{|E(C)\cap M|}
 \ge2^{\mathbb E_M|E(C)\cap M|}=2^{|E(C)|/3}.
\]
Average~\eqref{eq:cubic-matching-weight} and sum this inequality over the cycles.
This proves~\eqref{eq:cubic-length-weight}.
\end{proof}

The external input is J.~Edmonds, \emph{Maximum matching and a polyhedron with $0,1$-vertices},
Journal of Research of the National Bureau of Standards, Section~B \textbf{69B} (1965), 125--130,
\url{https://doi.org/10.6028/jres.069b.013}.
The original paper is available at
\url{https://nvlpubs.nist.gov/nistpubs/jres/69B/jresv69Bn1-2p125_A1b.pdf}.
The weighted inequalities do not by themselves establish the edge-half property.

\subsection{Exact formulas across a three-edge cut}

Let a nontrivial three-edge cut split a three-edge-connected simple cubic graph into connected shores $N_1,N_2$.
The three cut endpoints on each shore are distinct.
Otherwise an internal edge at a repeated endpoint, together with the remaining cut edge, gives a two-edge cut.
Label corresponding terminals by $1,2,3$.
Put $c_i=c(N_i)$, and let $p_{jk}$ and $q_{jk}$ count terminal paths in the two shores.
Then
\begin{equation}\label{eq:cubic-three-cut}
 S=\sum_{1\le j<k\le3}p_{jk}q_{jk},\qquad
 c(G)=c_1+c_2+S.
\end{equation}
The cut edge at terminal $j$ belongs to
\[
 \sum_{k\ne j}p_{jk}q_{jk}
\]
cycles, with subscripts interpreted as unordered pairs.
Hence some cut edge satisfies the edge-half bound whenever
\[
 S\le3(c_1+c_2).
\]
For an internal edge $e$ of $N_1$, let $p_{jk}(e)$ count terminal paths using $e$.
Its exact cycle frequency is
\begin{equation}\label{eq:cubic-three-cut-internal}
 c_e(G)=c_e(N_1)+\sum_{j<k}q_{jk}p_{jk}(e).
\end{equation}

One sufficient extension would be the following weighted three-terminal statement.
For every connected simple graph $N$ with three degree-two terminals and all other vertices cubic, and every $w_{jk}\ge0$, some edge $e$ would satisfy
\begin{equation}\label{eq:cubic-three-pole-candidate}
 2c_e(N)+2\sum_{j<k}w_{jk}p_{jk}(e)
 \le c(N)+1+\sum_{j<k}w_{jk}p_{jk}.
\end{equation}
This statement is not proved here.
If it holds, apply it with $w_{jk}=q_{jk}$.
The other shore contains a cycle, since it has $(3|V(N_2)|-3)/2$ edges and odd order at least three.
Thus $c_2\ge1$, and~\eqref{eq:cubic-three-cut-internal} would give an edge with frequency at most $c(G)/2$.

A simpler terminal-visitation bound is false.
Take the cubic graph with vertex set $\{0,\ldots,13\}$ and edge list
\[
\begin{split}
 &(0,5),(1,6),(2,7),(0,8),(1,8),(3,8),(0,9),\\
 &(2,9),(4,9),(2,10),(5,10),(6,10),(1,11),(4,11),\\
 &(7,11),(3,12),(4,12),(6,12),(3,13),(5,13),(7,13).
\end{split}
\]
It is three-edge-connected and has girth five.
After deleting vertex $2$, there are $50$ paths between terminals $7$ and $9$.
Exactly $26$ visit the third terminal $10$.
Thus even this restricted class does not guarantee that at most half the terminal paths visit the third terminal.
This does not disprove~\eqref{eq:cubic-three-pole-candidate}.

\subsection{A counterexample to the mean-length route}

A sufficient condition for the edge-half property would be mean cycle length at most $3n/4$.
This condition fails even at order eight.
Replace one vertex of $K_{3,3}$ by a triangle, attaching its three old incident edges to distinct triangle vertices.
The resulting simple cubic graph has cycle polynomial
\[
 z^3+3z^4+6z^5+6z^6+6z^7+6z^8.
\]
Indeed, $K_{3,3}$ has nine four-cycles and six six-cycles.
The replaced vertex belongs to six cycles of each length.
Each such cycle has two lifts, whose lengths increase by one and two.
The remaining three four-cycles persist, and the new triangle contributes one cycle.
The mean cycle length is therefore
\[
 \frac{171}{28}>6=\frac{3\cdot8}{4}.
\]
Nevertheless, every new triangle edge belongs to exactly $13$ of the $28$ cycles.
Thus the mean-length condition is stronger than the edge-half conclusion and cannot prove it in general.

\subsection{Independent finite checks}

The script \texttt{check\_cubic\_edge\_frequency.py} enumerates connected simple cubic graphs using nauty's \texttt{geng}.
It counts cycles as connected nonzero vectors of the binary cycle space.
For all $621$ such graphs through order fourteen, it verifies every triangle and induced-square identity above.
It also verifies~\eqref{eq:cubic-matching-weight} for all $8\,140$ perfect matchings in the graphs of order at least six.
The order-eight counterexample is checked separately, including its cycle polynomial and three triangle-edge frequencies.
The exact output is saved in \texttt{cubic-edge-frequency-results.json}.
These checks support the arguments but are not premises of the propositions.
